\textbf{Время:} 01:45

\textbf{Тезис:} зависимости фемтоскопических радиусов от среднего импульса пары ограничивают возможные модели расширения.

\textbf{На экране:} сначала горизонтальная ось k\_T и вертикальная ось радиуса в фемтометрах; затем легенда направлений и заполнение маркеров; после чтения кривых — выделенная фраза справа и вопрос о моделях расширения.

\textbf{Текст:} Теперь вместо корреляционной функции показаны радиусы, полученные из её формы. По горизонтали — k-T, модуль среднего поперечного импульса пары; по вертикали — радиус в фемтометрах. С ростом k-T мы переходим к парам с большим средним поперечным импульсом.

Сначала прочитаем легенду. Красные точки — out, синие — side, чёрные — long. Заполненные маркеры относятся к отрицательным пионам, незаполненные — к положительным. Все эти данные относятся к наиболее центральному классу 0–10 процентов. Сравнивать зависимость от импульса нужно внутри одной серии, сохраняя направление и заряд.

Общая тенденция — уменьшение радиусов с ростом k-T. Именно она сформулирована справа. Это согласуется с только что показанной картиной: разные импульсы выделяют разные области однородности. Речь не о том, что весь источник стал меньше во время одного и того же измерения.

Вопрос под выделенной фразой задаёт следующий уровень анализа: какие модели расширения воспроизводят эти зависимости? Подходящая модель должна описать несколько направлений и разные отборы, а не только одно значение радиуса. Один этот график ещё не выбирает единственное описание динамики.

\textbf{Переход:} В радиусе out присутствует и временная информация. Посмотрим, почему из него нельзя сразу получить длительность испускания.

\textbf{Справка:} А. Краева (STAR), Universe 10, 188 (2024), рис. 5, §3–4. Au+Au, √s\_NN=3 ГэВ, центральность 0–10\%. Это трёхмерные LCMS-радиусы; отдельного ограничения по быстроте пары в данном варианте анализа нет, но выборка ограничена регистрацией. Вертикальные отрезки и затенённые области показывают статистическую и систематическую неопределённость; их оценка не входит в основной рассказ. Точки и исходные обозначения сохранены. Зарядовые различия могут зависеть от остаточного электрического поля, распадов и других механизмов; один геометрический аргумент не определяет причину. Интерпретация различий зарядов через длительность требует дополнительных предположений.
